\corrigesChapitre{4}{Solutions, solubilité et pH}
\par\noindent\begin{minipage}{\linewidth}
\begin{corrige}{1}{Le vocabulaire}
\begin{questions}
\item L’eau est le solvant ; le sucre, le soluté.
\item Non : il contient plusieurs espèces, même s’il est homogène et transparent.
\end{questions}
\end{corrige}
\end{minipage}
\par\noindent\begin{minipage}{\linewidth}
\begin{corrige}{2}{Deux liquides}
\begin{questions}
\item Non dans les conditions usuelles de l’expérience.
\item Il est hétérogène : on distingue deux phases.
\end{questions}
\end{corrige}
\end{minipage}
\par\noindent\begin{minipage}{\linewidth}
\begin{corrige}{3}{Trois pH}
\begin{questions}
\item pH 3 : acide ; pH 7 : neutre ; pH 10 : basique.
\item Non : le pH est un nombre sans unité.
\end{questions}
\end{corrige}
\end{minipage}
\par\noindent\begin{minipage}{\linewidth}
\begin{corrige}{4}{Un filtre en papier}
\begin{questions}
\item Oui, si les grains sont retenus par le filtre.
\item Non : les espèces dissoutes traversent le filtre avec l’eau.
\end{questions}
\pcFigureEx{figures/cor-filtration.pdf}{Le filtre retient les grains de sable ; les espèces dissoutes traversent le filtre ordinaire avec l’eau.}
\end{corrige}
\end{minipage}
\par\noindent\begin{minipage}{\linewidth}
\begin{corrige}{5}{Une solution saturée}
\begin{questions}
\item La solubilité a été atteinte : la solution est saturée, avec du solide non dissous.
\item La solubilité peut varier avec la température.
\end{questions}
\end{corrige}
\end{minipage}
\par\noindent\begin{minipage}{\linewidth}
\begin{corrige}{6}{Adapter une quantité}
\begin{questions}
\item $20\times250/100=50\unite g$.
\item Non : $10\unite g$ restent non dissous à cette température, après équilibre.
\end{questions}
\par\smallskip\noindent\textbf{Pour comprendre.}\enspace À température et solvant identiques, la limite fournie donne $20\times250/100=50\unite g$. Sur les 60 g ajoutés, au plus 50 g se dissolvent ; un dépôt d’environ 10 g persiste dans ce modèle.\par
\end{corrige}
\end{minipage}
\par\noindent\begin{minipage}{\linewidth}
\begin{corrige}{7}{Conserver la masse}
\begin{questions}
\item $108\unite g$.
\item Non : la conservation de la masse n’impose pas l’additivité des volumes.
\end{questions}
\end{corrige}
\end{minipage}
\par\noindent\begin{minipage}{\linewidth}
\begin{corrige}{8}{Lire un encadrement}
\begin{questions}
\item La masse maximale est au moins $14\unite g$ et inférieure à $16\unite g$.
\item Elle est au moins $28\unite g$ et inférieure à $32\unite g$ pour $100\unite{mL}$, à la même température.
\end{questions}
\par\smallskip\noindent\textbf{Pour comprendre.}\enspace La masse maximale est au moins 14 g, mais strictement inférieure à 16 g dans les conditions de l’essai. Pour un volume d’eau double : au moins 28 g et moins de 32 g. Les essais ne donnent pas une valeur exacte à eux seuls.\par
\end{corrige}
\end{minipage}
\par\noindent\begin{minipage}{\linewidth}
\begin{corrige}{9}{Le gaz d’une boisson}
\begin{questions}
\item Une partie du gaz dissous s’échappe lorsque les conditions changent, notamment la pression.
\item Non : un dégazage peut être une transformation physique sans création d’une nouvelle espèce.
\end{questions}
\end{corrige}
\end{minipage}
\par\noindent\begin{minipage}{\linewidth}
\begin{corrige}{10}{Choisir une méthode}
\begin{questions}
\item Une distillation adaptée, réalisée en classe, peut recueillir l’eau après vaporisation puis liquéfaction.
\item L’évaporation de l’eau peut laisser le sel. Cela ne suffit pas à identifier exactement ce résidu.
\end{questions}
\pcFigureEx{figures/cor-distillation.pdf}{Vaporiser puis liquéfier pour recueillir l’eau ; une évaporation peut laisser le sel, sans recueillir l’eau.}
\end{corrige}
\end{minipage}
\par\noindent\begin{minipage}{\linewidth}
\begin{corrige}{11}{Interpréter une baisse de pH}
\begin{questions}
\item Non : 7,8 reste supérieur à 7 dans les conditions indiquées.
\item Oui : le pH a diminué. Acidification signifie une évolution vers un pH plus faible, pas forcément un passage sous 7.
\end{questions}
\end{corrige}
\end{minipage}
\par\noindent\begin{minipage}{\linewidth}
\begin{corrige}{12}{Le résultat trop précis}
\begin{questions}
\item Cette précision n’est pas justifiée par la méthode.
\item Une estimation entre 5 et 6, ou environ 5,5 si l’échelle et les conditions le permettent, avec la limite de lecture indiquée.
\end{questions}
\end{corrige}
\end{minipage}
\par\noindent\begin{minipage}{\linewidth}
\begin{corrige}{13}{Une comparaison équitable}
\begin{questions}
\item Non : la masse dissoute dépend aussi de la quantité de solvant et de la température.
\item Employer le même volume d’eau à la même température et laisser le temps d’atteindre une dissolution suffisante.
\end{questions}
\end{corrige}
\end{minipage}
\par\noindent\begin{minipage}{\linewidth}
\begin{corrige}{14}{Diluer sans deviner}
\begin{questions}
\item Il se rapproche de la neutralité, soit 7 vers 25 °C pour les solutions aqueuses considérées.
\item Non : le pH n’est pas proportionnel au volume de dilution. La valeur exacte exige des informations supplémentaires.
\end{questions}
\end{corrige}
\end{minipage}
\par\noindent\begin{minipage}{\linewidth}
\begin{corrige}{15}{L’eau prétendument potable}
\begin{questions}
\item Il manque des contrôles microbiologiques et chimiques adaptés ; l’aspect ne révèle pas tous les contaminants.
\item La filtration a pu retirer certaines particules. On ne peut conclure à la potabilité à partir de cette seule observation.
\end{questions}
\end{corrige}
\end{minipage}
\par\noindent\begin{minipage}{\linewidth}
\phantomsection\label{cor-defi-4}\pcSolutionDefi{4}{Préparer sans gaspiller}\begin{questions}
\item $30\times80/100=24\unite g$ se dissolvent au maximum ; il reste $4\unite g$.
\item $V=28\times100/30\approx93{,}3\unite{mL}$ d’eau, au total, à la température donnée.
\item Les données sont arrondies, la température varie et la dissolution prend du temps. Un petit excès raisonné de solvant peut être nécessaire, sans déverser les déchets sans consigne.
\end{questions}
\end{minipage}
\par\noindent\begin{minipage}{\linewidth}
\phantomsection\label{cor-activite-4}\pcSousTitre{Pistes pour l’investigation}
La limite est encadrée entre la dernière masse complètement dissoute et la première laissant un dépôt persistant. L’intervalle dépend de la taille des ajouts et de la durée d’agitation. La température et le volume d’eau doivent rester comparables. Pour le pH, prélever un petit échantillon sans remettre la bandelette dans le flacon initial et annoncer une estimation compatible avec l’échelle.\par\end{minipage}
